\documentclass[11pt]{article}
\usepackage[T1]{fontenc}
\usepackage{lmodern,amsmath,amssymb,amsthm}
\usepackage[margin=1in]{geometry}
\usepackage[hidelinks]{hyperref}
\newtheorem{theorem}{Theorem}
\newtheorem{corollary}{Corollary}
\newcommand{\Q}{\mathbb Q}
\newcommand{\F}{\mathbb F}
\title{Ranks of Two Elliptic K3 Surfaces Obtained by Quadratic Base Change}
\author{\small \href{mailto:admin@cosmosgeek.org}{admin@cosmosgeek.org}\qquad \href{mailto:b321014@iiit-bh.ac.in}{b321014@iiit-bh.ac.in}}
\date{Working draft for specialist review}
\begin{document}
\maketitle
\begin{abstract}
We compute the Mordell--Weil groups up to rank for two explicit K3 surfaces
arising from the rational elliptic surface $y^2=x^3-28x+t^2-48$.
Each has geometric rank four, and exactly three of these directions are
defined over the rational function field with rational constants.
Counting fibres in two finite extensions of a single residue field
gives the Picard-rank upper bound in each case. Over the compositum of
the two quadratic function fields the rational rank is at least four,
and a positive-rank parameter curve supplies infinitely many rational
specializations with four independent points.
\end{abstract}

\section{Surfaces and sections}
Let $E_t/\Q(t)$ denote the elliptic curve in the abstract. For
$d=48,33$ write
\[
 t=(s+d/s)/2,\qquad v=(s-d/s)/2,\qquad t^2-v^2=d,
\]
and let $S_d$ be the relatively minimal elliptic surface with this
generic fibre over the $s$-line.

\begin{theorem}
For $d=48,33$ the geometric Picard number of $S_d$ is $18$, and
\[
 \operatorname{rank} E_t(\overline{\Q}(s))=4,\qquad
 \operatorname{rank} E_t(\Q(s))=3 .
\]
The fourth direction is defined over $\Q(\sqrt7)$ for $d=48$ and
over $\Q(\sqrt{37})$ for $d=33$. The two surfaces have good
reduction at $31$ and $41$, respectively. Over the corresponding
finite field function fields the arithmetic and analytic ranks
are both four.
\end{theorem}

The original surface is rational: the coordinates $a=y-t$ and $x$
give a birational parametrization since $a(y+t)=x^3-28x-48$.
The discriminant is $16(87808-27(t^2-48)^2)$, with four simple
zeros. Its finite singular fibres have type $I_1$ and its fibre at
infinity has type $IV^*$. Thus Shioda--Tate gives geometric rank
$10-2-6=2$ over the original $t$-line \cite{SS}.

Write $P=(-2,t)$ and $Q=(-4,t)$. To exhibit two more directions,
set $x_0=0$ for $d=48$ and $x_0=-5$ for $d=33$. The roots
of $x^3-28x+d-48$ are $x_0$ and $x_\pm$, where
\[
 (x_+,x_-)=(2\sqrt7,-2\sqrt7)
 \quad\hbox{or}\quad ((5+\sqrt{37})/2,(5-\sqrt{37})/2)
\]
respectively. Put $R=(x_0,v)$ and $D=(x_+,v)-(x_-,v)$ using the
elliptic group law. Under $s\mapsto d/s$, $P,Q$ are invariant
and $R,D$ anti-invariant. The conjugation on the constant quadratic
field fixes $P,Q,R$ and sends $D$ to $-D$.

Here are finite checks for independence. At $t=7$ the specialized
curve $E_7$ has equation $y^2=x^3-28x+1$. Taking $s=8$ for
$d=48$ yields $v=1$. For $P,Q,R$ the images after reduction and
multiplication by $\#E_7(\F_\ell)/5$ have the following discrete
logarithms in a cyclic group of order five:
\[
\begin{array}{c|r|c|ccc}
\ell & \#E_7(\F_\ell)&\text{generator}&P&Q&R\\\hline
19&25&(3,18)&1&2&1\\
71&75&(59,10)&0&1&2\\
83&90&(63,23)&1&0&3
\end{array}
\]
These entries may be verified by the chord and tangent law in
$\F_\ell$; the orders are obtained by counting $1+\chi(x^3-28x+1)$
and adding the point at infinity. The determinant of the three rows
is $6$. The specialization has no rational $5$-torsion, since
$\#E_7(\F_7)=12$. Thus the three points, and in particular $P,Q$,
are independent. Indeed any relation up to torsion, chosen primitive
after passing to the free quotient, would give a relation modulo five.
The reduction at $19$ also distinguishes $(x_+,1)$ and $(x_-,1)$
at a prime of $\Q(\sqrt7)$ above $19$. Geometric torsion on the
surface injects into the component group of a $IV^*$ fibre in
characteristic zero: the formal group and the additive identity
component have trivial torsion. A torsion $D$ would consequently have
order dividing three, whereas its image at $19$ is nonzero in a
group of order $25$.

For $d=33$ use $s=11$, so $t=7,v=4$. Both $11$ and $47$ split in
$\Q(\sqrt{37})$; the good reductions have orders $11$ and $52$.
The prime-to-residue-characteristic torsion injection at these two
primes excludes every torsion prime. Thus both $(-5,4)$ and the
difference of the distinct specialized conjugate points are nontorsion.
Together with the two involutions this proves independence of all four
sections over the geometric constant field in both cases.

\section{Fibres and specialization}
The double cover branches over $t^2=d$, away from the singular fibres.
It gives eight $I_1$ fibres and two $IV^*$ fibres, at zero and infinity.
The elliptic surface with section over $\mathbb P^1$ has
$\chi(\mathcal O)=24/12=2$ and is a K3 surface by the canonical
bundle formula \cite{SS}. The trivial lattice has rank $14$.

We give some detail on good reduction, since a count of smooth
specializations of the generic curve alone would not justify the
Picard bound. In the two residue characteristics in the theorem,
$27(t^2-48)^2-87808$ has simple roots and is nonzero at $t^2=d$.
This follows by differentiating and testing $t=0$, $t^2=48$ and
$t^2=d$ modulo the stated prime. The nodal fibres there have smooth
total space. The integral minimal equation at each of the two ends
has the form
\[
 y^2=x^3+z^4(Ax+B(z)),\qquad A=-28,\quad B(0)=b^2\ne0,
\]
with $b=d/2$ or $1/2$, both units. The same resolution works over
the local integers and their residue field. After blowing up the
origin, its $z$-chart is $Y^2=zX^3+z^3AX+z^2B(z)$. The other
nonempty exceptional chart is smooth, and only its $z$-chart origin
requires a second blowup. The latter's $z$-chart has $V=\pm b$
on its exceptional divisor. In the $X$-chart the equation is
\[
 V^2=Z^2B(XZ)+X^2 Z(1+AZ^2).
\]
Only its origin is singular. In the completed local chart factor the
difference of squares using the unit square root of $B(XZ)$. With
$U_\pm=V\pm Z\sqrt{B(XZ)}$ it becomes
$U_+U_-=X^2 h(U_+-U_-)$ for a unit $h$.
Replacing the two coordinates by $U_++X^2h,U_--X^2h$
gives product $-X^4h^2$. This is a split $A_3$ singularity
formally equivalent to $uv=X^4$. Successive blowups of the singular
sections resolve it in both characteristics. The coordinate changes
have invertible linear part and all denominators used are units.
This supplies compatible smooth proper models. The fibre components
are split (equivalently the leading coefficient $B(0)$ is a unit
square in Tate's $IV^*$ test).

Specialization of divisor classes in a smooth proper family is
injective on the geometric N\'eron--Severi group tensored with
$\Q$, by the cycle class injection and smooth proper base change.
Thus the trivial lattice and independent section classes give
eighteen independent directions also after reduction. All these
directions are defined over the stated residue field since
$(7/31)=(37/41)=1$.

\section{A finite Frobenius certificate}
For $q=p^j$ in the table below and $u\ne0$ in $\F_q$, put
$t_u=(u+d/u)/2$ and
\[
 S_q=-\!\sum_{u\in\F_q^\times}\sum_{x\in\F_q}
       \chi_q(x^3-28x+t_u^2-48),\qquad \chi_q(0)=0 .
\]
The raw character sum counts a projective nodal fibre correctly as
well. Each split $IV^*$ fibre has $7q+1$ points: its seven rational
components form a tree. Therefore $\#S_d(\F_q)=q^2+1+14q-S_q$.
The trace formula for a K3 surface gives the trace $14q-S_q$ on
$H^2$. Subtracting the eighteen known directions gives trace
$-S_q-4q$ on their four-dimensional orthogonal complement $W$.
The trivial lattice together with the nondegenerate Mordell--Weil
height pairing makes the known subspace nondegenerate.

\[
\begin{array}{r|r|r|r|r|r}
d&p&S_p&S_{p^2}&\operatorname{tr}(W)&\operatorname{tr}(\mathrm{Fr}_p^2|W)\\\hline
48&31&-124&-3720&0&-124\\
33&41&-184&-5320&20&-1404
\end{array}
\]
Here is reproducible pseudocode. For $j=2$ use
$\F_p[z]/(z^2-3)$; $3$ is a nonsquare at both primes.
\begin{verbatim}
K := field with p^j elements as specified
chi := array indexed by K, initially -1
for x in K: chi[x*x] := 1
chi[0] := 0
H,T := zero integer arrays indexed by K
for x in K: H[x^3 - 28*x] += 1
for u in K, u != 0: T[((u + d/u)/2)^2 - 48] += 1
S := -sum(H[a]*T[b]*chi[a+b] for a,b in K)
\end{verbatim}

By Poincar\'e duality, if the constant term of the monic Frobenius
polynomial of $W$ is $-p^4$, the middle coefficient must vanish.
The middle coefficients given by Newton's identities are $62$ and
$902$, respectively. The reciprocal polynomials are therefore
\[
 X^4+62X^2+31^4,\qquad
 X^4-20X^3+902X^2-33620X+41^4 .
\]
A further geometric divisor would supply an eigenvalue $p\zeta$
with $\zeta$ a root of unity. In the first polynomial this gives
$\zeta^2+\zeta^{-2}=-2/31$, which is not an algebraic integer.
In the second, $h=\zeta+\zeta^{-1}$ satisfies
$1681h^2-820h-2460=0$. If $h$ has degree two its trace is $20/41$;
if rational it must be one of $-2,-1,0,1,2$, none a root.
This proves the Picard upper bound over both geometric residue fields
and hence in characteristic zero. Shioda--Tate now gives geometric
rank four. The four exhibited directions span over $\Q$, with
constant-field conjugation negating exactly $D$, proving the rational
rank statement. Over each residue field all four directions are
rational. The Tate theorem for elliptic K3 surfaces and its
rank equivalence with BSD for the generic fibre over the finite
function field give analytic rank four there \cite{ASD,Ulmer}.

\section{Four independent points after both base changes}
Over $\Q(t,\sqrt{t^2-48},\sqrt{t^2-33})$, the four rational
sections with $x=-2,-4,0,-5$ are independent: the two radical signs
are distinct characters over $\Q(t)$, and the first two sections
remain independent. This conclusion is a lower bound. The remaining
quadratic-character component, associated to $(t^2-48)(t^2-33)$,
has not been computed.

For a concrete specialization let $t=817/56$,
$v=719/56$, and $w=751/56$, so $v^2=t^2-48$, $w^2=t^2-33$.
On the specialized curve the four sections have coordinates
$(-2,t),(-4,t),(0,v),(-5,w)$. The rows below give their multiples by
$\#E(\F_\ell)/5$, expressed relative to the first nonzero image:
\[
\begin{array}{r|r|rrrr}
\ell&\#E(\F_\ell)&-2&-4&0&-5\\\hline
103&105&1&4&2&2\\
197&190&1&2&4&1\\
227&235&1&2&4&3\\
257&280&1&0&0&4
\end{array}
\]
The determinant is $-24$. The reduction at $13$ has order sixteen, so
there is no rational $5$-torsion. The table proves rank at least four.

Finally, writing $u=t-v$ makes the second square condition
$Y^2=u^4-36u^2+2304$, $Y=2uw$. This curve is birational to
\[
 B:\ Z^2=(X-36)(X-96)(X+96),\qquad
 X=2(Y+u^2),\quad Z=2u(X-36).
\]
The inverse has $u=Z/(2(X-36))$ and $Y=X/2-u^2$.
The point $T=(256,3520)$ has double $(100,224)$ giving $t=817/56$.
At good primes $7$ and $13$, $B$ has respectively eight and twelve
points and $T$ reduces to points of orders four and three. If
torsion, its order would divide four by prime-to-residue-characteristic
injectivity, a contradiction. The multiples $(2+kM)T$ with $M$ a
multiple of the orders of $B$ modulo the four tabulated primes and
modulo $13$ reproduce the reductions of $2T$ at these good places.
The inverse formula is integral there at $2T$: its denominators
are units at the five listed odd primes.
Exclude finitely many poles. The table and torsion exclusion persist.
The maps to $t$ and to $j(E_t)$ are nonconstant, so infinitely many
distinct $j$-invariants occur.

\begin{thebibliography}{9}
\bibitem{SS} M. Sch\"utt and T. Shioda,
\emph{Elliptic surfaces}, \href{https://arxiv.org/abs/0907.0298}{arXiv:0907.0298}.
\bibitem{ASD} M. Artin and H. P. F. Swinnerton-Dyer,
\emph{The Shafarevich--Tate conjecture for pencils of elliptic curves on K3 surfaces},
Invent. Math. \textbf{20} (1973), 249--266.
\bibitem{Ulmer} D. Ulmer, \emph{Curves and Jacobians over function fields},
\href{https://math.arizona.edu/~ulmer/research/papers/2014b.pdf}{author's manuscript}.
\end{thebibliography}
\end{document}
